% Part: intuitionistic-logic
% Chapter: tableaux
% Section: rules

\documentclass[../../../include/open-logic-section]{subfiles}

\begin{document}

\olfileid{int}{tab}{rul}

\olsection{直觉主义逻辑的规则}

联结词 $\land$ 和 $\lor$ 的规则与标准的带号命题分析表规则相同，只是加上了前缀。
在每种情况下，对带号公式
$\sFmla{S}{!A}[\sigma]$ 应用规则所产生的新公式也都带有前缀~$\sigma$。
这在直观上应该很清楚：例如，如果 $!A
\land !B$ 在（由~$\sigma$ 命名的世界中）为真，那么 $!A$ 和 $!B$
在~$\sigma$ 处为真（而不是在任何其他世界为真）。我们将
$\land$ 和 $\lor$ 的规则汇总于\olref{tab:prop-rules}。

\begin{table}
  \[\def\arraystretch{3}\begin{array}{|c|c|}
    \hline
    \AxiomC{\sFmla{\True}{!A \land !B}[\sigma]}
    \RightLabel{\TRule{\True}{\land}}
    \UnaryInfC{\sFmla{\True}{!A}[\sigma]}
    \noLine
    \UnaryInfC{\sFmla{\True}{!B}[\sigma]}
    \DisplayProof
    &
    \AxiomC{\sFmla{\False}{!A \land !B}[\sigma]}
    \RightLabel{\TRule{\False}{\land}}
    \UnaryInfC{$\sFmla{\False}{!A}[\sigma] \quad \mid \quad
      \sFmla{\False}{!B}[\sigma]$}
    \DisplayProof
    \\[2ex]
    \hline
    \AxiomC{\sFmla{\True}{!A \lor !B}[\sigma]}
    \RightLabel{\TRule{\True}{\lor}}
    \UnaryInfC{$\sFmla{\True}{!A}[\sigma] \quad \mid \quad
      \sFmla{\True}{!B}[\sigma]$}
    \DisplayProof
    &
    \AxiomC{\sFmla{\False}{!A \lor !B}[\sigma]}
    \RightLabel{\TRule{\False}{\lor}}
    \UnaryInfC{\sFmla{\False}{!A}[\sigma]}
    \noLine
    \UnaryInfC{\sFmla{\False}{!B}[\sigma]}
    \DisplayProof
    \\[2ex]
    \hline
  \end{array}\]
  \caption{$\land$ 和 $\lor$ 的带前缀分析表规则}
  \ollabel{tab:prop-rules}
\end{table}

闭包条件与普通分析表类似，不过我们要求不仅公式必须匹配，前缀也必须匹配。
实际上，我们可以放宽一些：由于在直觉主义模型中，公式一旦为真，则保持为真，
因此 $!A$ 不可能在~$\sigma$ 处为真，却在某个可达的前缀~$\sigma.{*}$ 处为假。
所以，如果某个分支对于某个前缀 $\sigma$ 和公式~$!A$，同时包含
\[
\sFmla{\True}{!A}[\sigma] \quad\text{和}\quad \sFmla{\False}{!A}[\sigma.{*}]
\]
则该分支闭合。注意，如果符号相反，即如果该分支包含
\[
\sFmla{\False}{!A}[\sigma] \quad\text{和}\quad \sFmla{\True}{!A}[\sigma.{*}]
\]
则仅当 $*$ 是空序列时该分支才闭合。

此外，如果一个分支包含~$\sFmla{\True}{\bot}[\sigma]$，则该分支闭合。

建立假设的规则也与普通分析表一样，只不过假设总是使用
前缀~$1$。（具体使用哪个前缀并不重要，只要所有假设使用相同的前缀即可。）
因此，例如，我们说
\[
!B_1, \dots, !B_n \Proves !A
\]
当且仅当存在一个从假设
\[
\sFmla{\True}{!B_1}[1], \dots, \sFmla{\True}{!B_n}[1],
\sFmla{\False}{!A}[1].
\]
出发的闭合分析表。

对于条件句~$\lif$，其规则不同于经典情况和模态情况。
$\TRule{\lif}{\True}$ 规则通过在两个不同的分支上分别添加 $\sFmla{\True}{!A}[\sigma.{*}]$ 和 $\sFmla{\False}{!B}[\sigma.{*}]$，
来扩展包含 $\sFmla{\True}{!A \lif !B}[\sigma]$ 的分支。
它只能针对在应用该规则的分支上\emph{已经}出现的前缀~$\sigma.{*}$ 来应用。
我们称这样的前缀（在该分支上）“已使用的”。（由于 $\sigma.{*}$
包含 $\sigma$ 本身，该规则总可以通过在各自独立的分支上添加带前缀的带号公式
$\sFmla{\True}{!A}[\sigma]$ 和
$\sFmla{\False}{!B}[\sigma]$ 来应用。）

$\TRule{\lif}{\False}$ 规则通过在同一分支上同时添加 $\sFmla{\True}{!A}[\sigma.n]$ 和 $\sFmla{\False}{!B}[\sigma.n]$，
来扩展包含 $\sFmla{\False}{!A \lif !B}[\sigma]$ 的分支，
其中 $\sigma.n$ 是该分支上一个新的前缀。

$\lnot$ 的规则可类似地定义（利用 $\lnot !A$ 定义为 $!A \lif \lfalse$）。

这些规则如\olref{tab:rules-lif-lnot}所示。

\begin{table}
  \[\def\arraystretch{3}\begin{array}{|c|c|}
    \hline
    \AxiomC{\sFmla{\True}{\lnot !A}[\sigma]}
    \RightLabel{\TRule{\True}{\lnot}}
    \UnaryInfC{$\sFmla{\False}{!A}[\sigma.{*}]$}
    \DisplayProof
    &
    \AxiomC{\sFmla{\False}{\lnot !A}[\sigma]}
    \RightLabel{\TRule{\False}{\lnot}}
    \UnaryInfC{\sFmla{\True}{!A}[\sigma.n]}
    \DisplayProof
    \\[1ex]
    \text{（$\sigma.{*}$已使用）} & \text{（$\sigma.n$是新的）}\\
    \hline
    \AxiomC{\sFmla{\True}{!A \lif !B}[\sigma]}
    \RightLabel{\TRule{\True}{\lif}}
    \UnaryInfC{$\sFmla{\False}{!A}[\sigma.{*}] \quad \mid \quad
      \sFmla{\True}{!B}[\sigma.{*}]$}
    \DisplayProof
    &
    \AxiomC{\sFmla{\False}{!A \lif !B}[\sigma]}
    \RightLabel{\TRule{\False}{\lif}}
    \UnaryInfC{\sFmla{\True}{!A}[\sigma.n]}
    \noLine
    \UnaryInfC{\sFmla{\False}{!B}[\sigma.n]}
    \DisplayProof
    \\[1ex]
    \text{（$\sigma.{*}$已使用）} & \text{（$\sigma.n$是新的）}\\
    \hline
  \end{array}\]
  \caption{$\lnot$ 和 $\lif$ 的带前缀分析表规则}
  \ollabel{tab:rules-lif-lnot}
\end{table}

\end{document}